\paragraph{Schedules and warmup.} Cosine annealing was proposed together with warm restarts for SGD in~\cite{loshchilov2016sgdr}; we use its single-cycle form without restarts, as is common for language models. Warmup in the large-batch regime is motivated in~\cite{goyal2017accurate}, and for Adam by the variance of the adaptive learning rate in the first steps~\cite{liu2019variance}. For GPT-style models,~\cite{kalra2024why} study the mechanisms by which warmup permits larger peak learning rates, and~\cite{kosson2024analyzing} analyse when warmup can be reduced or removed. Our warmup ablation is a small-scale probe of the same phenomenon in the specific setting of cross-schedule comparison, where the presence of warmup is a confound.

\paragraph{Decay phases.} The cooldown (decay) stage of warmup-stable-decay schedules has been analysed in~\cite{dremov2025training}. Our step-decay and cosine schedules both contain a decay phase, but we do not study the dynamics of that phase; we only compare final loss.

\paragraph{Schedule-free optimisation.} Schedule-Free SGD/AdamW~\cite{defazio2024road} interpolates between an averaged iterate and a base iterate and is evaluated at the average, removing the need to choose a decay horizon. We reimplement the AdamW variant from the algorithm description (not from the authors' code), including our choice of warmup; differences from the reference implementation may exist.

\paragraph{Benchmarks.} Large-scale benchmarks~\cite{schmidt2020descending,dahl2023benchmarking} stress equal tuning budgets and show that optimiser rankings are sensitive to the protocol. Our study asks a complementary, narrower question: sensitivity to a single hyperparameter, the peak learning rate, at a tiny scale.
